\section{An explicit equation in ten unrestricted integers}
\label{sec:signed10}

The smallest unknown count in this package is achieved by the following
fully fixed equation:
\[
 F_{10}(f,g,h,k,l,w,x,y,z,m)=0,
 \qquad(f,g,h,k,l,w,x,y,z,m)\in\mathbb Z^{10}.
\]
Its division-free circuit uses \textbf{1,442 binary arithmetic operations}
and has exact total degree
\[
 126496\cdot5^{31}+27064
 =589042901992797851562527064.
\]
The general ten-integer reduction is due to Sun's 2026 preprint; this
report does not claim a new universal unknown-count theorem.  Our
contribution here is its completely specified Con(ZFC) specialization,
together with three local reductions of its degree and arithmetic circuit.
The two primary sources are Sun's 2021 paper, especially the explicit
proof of Theorem~3.1, and his 2026 preprint, Sections~3, 6 and~7.\footnote{
Z.-W.~Sun, \emph{Further results on Hilbert's Tenth Problem},
Sci. China Math. 64 (2021), 281--306,
\url{https://arxiv.org/html/1704.03504v7};
Z.-W.~Sun, \emph{Ten unknowns for Hilbert's tenth problem over the integers},
2026 preprint, \url{https://arxiv.org/html/2609.01594v2}.}
The literal Sun final formula is retained as a reference variant.

\subsection{Fixing every source coefficient}

Let $r_1,\ldots,r_{22}$ denote the quadratic residuals of the positive
31-variable system, with the complete new four-port tuple.  Define
\[
 P_0(n_1,\ldots,n_{31})
 =\sum_{j=1}^{22}r_j(n_1+1,\ldots,n_{31}+1)^2.
\]
The preceding reduction proves that a nonnegative zero of $P_0$ exists
exactly when ZFC is inconsistent.  Its degree is four, and $P_0(0)>0$:
the fifth residual $E-X(a-E)$ equals one at the all-one parent assignment.
Apply Sun's source construction to
\[
 P(a,n_1,\ldots,n_{31})=P_0(n_1,\ldots,n_{31}),\qquad a=0.
\]
The input $a$ is dummy.  This direct specialization avoids any unspecified
index for a recursively enumerable set.

For $\alpha\in\mathbb N^{32}$ with $|\alpha|\le4$, write $p_\alpha$ for
the corresponding ordinary coefficient of $P$ and set
\[
 c_\alpha=(4-|\alpha|)!\prod_{j=0}^{31}\alpha_j!\,p_\alpha,
 \qquad e(\alpha)=\sum_{j=0}^{31}\alpha_j5^j.
\]
Exactly 402 nonzero weighted coefficient recipes are retained in
\texttt{signed10/data/constants.json}; unlisted coefficients are zero.
The verifier reconstructs all of them from the shifted residual squares.
All nonzero records have $\alpha_0=0$.  The checked coefficient bound is
\[
 |c_\alpha|+1<2^{4269855}.
\]
Fix
\[
 \beta_S=2^{4269882}=2^{2\cdot2134941}.
\]
Because $48\cdot33^4=56924208<2^{27}$ and $|p_\alpha|\le|c_\alpha|$,
\[
 \beta_S>33^4\,4!\,2\max_\alpha|p_\alpha|.
\]
This is precisely the source bound in Sun 2021 for $\nu=31$, $\delta=4$
and $p=2$.  The additional constant is supplied both by its exact power
and by a compressed binary payload.  The binding of every fixed leaf is
\texttt{signed10/data/INSTANCE.json}.

\subsection{The complete source polynomials}

Put $t=5^{31}$ and $L_c=5^{32}=5t$.  Define
\begin{align*}
 b&=1+3f,&\mathcal B&=\beta_S b^4,\\
 \ell&=\sum_{j=1}^{31}\mathcal B^{5^j},&
 D_c&=\sum_{|\alpha|\le4}c_\alpha\mathcal B^{L_c-e(\alpha)},\\
 N_0&=\mathcal B^{t+1},&N_1&=4\mathcal B^{9t+1},\\
 N&=N_0N_1,&M_c&=N_0-1-(b-1)\ell,\\
 T_c&=M_c+(\mathcal B-2)\mathcal B^{L_c}N_0,&
 J_c&=2(1+g)^4D_c+\sum_{i=0}^{9t}\mathcal B^{i+1},\\
 S_c&=g+J_cN_0,&R_c&=(S_c+T_c+1)N+T_c+1,\\
 X&=(N-1)R_c,&Y&=N^2,\\
 \mathscr C&=b\mathcal B^t.&&
\end{align*}
All powers and summation limits are fixed integers.  The apparent division
by $p-1$ in Sun's general formula disappears at $p=2$.  The mask $M_c$
is his digit mask in telescoped form.  The bound has the prescribed form
$\mathscr C=2^{4269882t}b^{4t+1}$.

Sun's source theorem supplies, from a nonnegative zero of $P_0$, witnesses
with $f>0$, $b$ a square and a power of two, $b\le g<\mathscr C$, and
$Y\mid\binom{2X}{X}$.  Conversely, $f\ne0$, $0\le g<2\mathscr C$, the
same square and power conditions on $b$, and that divisibility imply a
nonnegative zero of $P_0$.  The construction also gives the strict size
bounds $X>3b$ and $Y>\max\{b,256\}$ required by Sun 2026.  Thus all source
hypotheses are checked for this particular, fully fixed polynomial.

\subsection{The bridge and the shared Lucas witnesses}

Use the following polynomial abbreviations, in the displayed order:
\begin{align*}
 L&=lY,& U&=2LX,& V&=4gwY,& W&=bw,\\
 A&=U(V+1),& B&=2X+1,& C&=B+(A-2)h,\\
 P_n&=U^2V,& K&=X+1+k(P_n-2),\\
 S&=2A-5,&T&=3WC-2(W^2-1),\\
 R_2&=(P_n^2-4)K^2+4.&&
\end{align*}
Then set
\begin{align*}
 \Delta&=8\mathscr C^3gK^2
 -g^2\bigl(32(C-KL)^2\mathscr C^3+g^2K^2\bigr),\\
 O&=f^2l^2\Delta,\qquad \xi=zx,\\
 D&=(A^2-4)C^2+4,\qquad E=C^2D\xi,\\
 F&=4(A^2-4)E^2+1,\\
 G^*&=2(A+2)(A-2)^2E^2-1,\qquad \Omega=16E,\\
 H^*&=C+(B-C)F+\Omega Fy,\\
 I^*&=((G^*)^2-1)(H^*)^2+1.
\end{align*}
These are Sun 2026's formulas (6.1)--(6.7), followed by his modified Lucas
certificate (3.1)--(3.7).  The symbol $P_n$ only factors a repeated product.
The integer $S$ is always nonzero and odd.

His positive bridge supplies $b$ square, $R_2$ square, $S\mid T$,
$C=u_B(A)$, and $O>0$.  In fact $O\ge4$: on this branch $g\ge b\ge4$,
and $\Delta>0$ is a multiple of $g$; also $f,l\ne0$.
His reverse bridge states that
\[
 b\text{ square},\quad O>0,\quad\xi\ne0,\quad
 DFI^*\text{ square},\quad R_2\text{ square},\quad S\mid T
\]
imply a nonnegative zero of $P_0$.  These precise bridge statements are
the external Lucas-theoretic inputs used below.

\subsection{Three local simplifications}

\paragraph{A quintic merger.}
Use
\[
 \pi(v)=(v^3-v-1)(v^2+23),\qquad
 \Sigma=\xi S^5,
\]
and
\[
 r=T+Sy,\qquad
 \Theta=(r^3-rS^2-S^3)(r^2+23S^2).
\]
The polynomial $\pi$ has a root modulo every positive integer.  For an
odd prime other than 23, either its cubic factor has a simple root, or
that cubic is irreducible and its discriminant $-23$ is a square: the
Frobenius permutation is an even 3-cycle.  In the latter case the quadratic
factor has a simple root.  At 23 the cubic has the simple root 3.  At 2,
$-23\equiv1\pmod8$ has a 2-adic square root.  Hensel lifting and the Chinese
remainder theorem finish the argument.  Neither factor has a rational root.

If $S\mid T$ and $\xi\ne0$, prescribe a residue class of $y$ modulo $\xi$
so that $\Sigma\mid\Theta$; Sun's full residue interpolation for the Lucas
witness allows exactly this prescription.  Conversely, if a prime has
$\nu_p(S)>\nu_p(T)$, the two factors of $\Theta$ have valuations
$3\nu_p(T)$ and $2\nu_p(T)$, contradicting $S^5\mid\Theta$.  Hence
$\Sigma\mid\Theta$ forces $S\mid T$.  Moreover $\Theta\ne0$ whenever
$S\ne0$, by the absence of rational roots.  This replaces Sun's degree-six
merger with a degree-five one.

\paragraph{One square in place of two.}
Define
\[
 r_0=b\bigl(1+(3O-4)b^2z^2\bigr),\qquad
 r_1=R_2,\qquad r_2=DFI^*.
\]
For $z\ne0$, a square $r_0$ forces $b$ to be a positive square and
$O>0$.  Indeed its two factors are coprime.  The congruence
$b\equiv1\pmod3$ excludes zero and negative squares.  A zero factor would
force $b=1$, hence $f=O=0$ and $1-4z^2=0$, which is impossible.
The remaining factor is a
square of the form $1+(3O-4)(bz)^2$ with $bz\ne0$, so the standard Pell
positivity lemma gives $O>0$.

For completeness, the positive bridge has $O\ge4$, so $d=3O-4>0$ is
nonsquare because it is 2 modulo 3.  A positive Pell solution
$u^2-dv^2=1$ can be chosen with $b\mid v$: raise a fundamental Pell unit
to the order of its image modulo $b$.  Taking $z=v/b>0$ makes $r_0$ a
square.  Sun's interpolation theorem then chooses $\xi$ divisible by
this $z$, so $x=\xi/z$ is still an integer.  No witness is added.

\paragraph{Smaller separating weights.}
For the actual polynomial $r_0$, $z\ne0$ implies $|r_0|>1$.  If $r_0=0$
or $r_0=\pm1$, its factorization forces $b=1$, hence $f=O=0$; but then
$r_0=1-4z^2\le-3$, a contradiction.  Also
\[
 \sqrt{|r_1|}\le P_n^2+K^2+2,
\]
because $|r_1|\le(P_n^2+4)K^2+4$ and the difference from the square of
the right-hand side is
$P_n^4+K^4+P_n^2K^2+4P_n^2\ge0$.
Thus the possibly signed weights
\[
 w_0=1,\qquad w_1=r_0,\qquad w_2=r_0(P_n^2+K^2+3)
\]
satisfy
\[
 |w_1|>\sqrt{|r_0|},\qquad
 |w_2|>\sqrt{|r_0|}+|w_1|\sqrt{|r_1|}.
\]
If a signed sum of these weighted roots is rational, any Galois
automorphism changing a root has a largest changed index.  Its term has
magnitude at least the corresponding weight and cannot be cancelled by
the earlier terms.  Hence every root is rational, and the integer
radicands are squares.  The same argument covers negative radicands by
complex absolute values, and zero radicands contribute no term.

\subsection{The final polynomial, without radicals or division}

Set
\[
 X_n=\Sigma m-\Theta,\qquad
 A_i=\Sigma^2w_i^2r_i\quad(0\le i\le2),
\]
and then
\[
 V_n=X_n^2+A_2-A_0-A_1,\qquad
 W_n=V_n^2+4X_n^2A_2-4A_0A_1.
\]
The requested equation is
\begin{equation}
 \boxed{F_{10}=W_n^2-16V_n^2X_n^2A_2=0.}
 \label{eq:teninteger}
\end{equation}
All preceding letters are polynomial abbreviations, not additional
unknowns.  Its formal norm identity is
\[
 F_{10}=\prod_{\varepsilon\in\{\pm1\}^3}
 \left(X_n+\sum_{i=0}^2\varepsilon_i\sqrt{A_i}\right).
\]
This identity explains the construction; equation~\eqref{eq:teninteger}
itself contains no radicals.  Its 35-term formal expansion is included
in the archive and checked exactly against the eight signed factors.

For a nonnegative zero of $P_0$, the positive bridge, the merged-square
Pell construction and the full residue-interpolation theorem supply
$r_0,r_1,r_2$ square, $\Sigma\ne0$ and $\Sigma\mid\Theta$.  Choose the
integer $m$ to make a norm factor vanish.  Conversely, if $F_{10}=0$,
then $\Sigma=0$ would give $F_{10}=\Theta^8\ne0$.  Thus $\xi=zx\ne0$
and $z\ne0$, so the structural separating-weight argument applies.
A vanishing norm factor forces the three radicands to be squares and
$\Sigma\mid\Theta$.  The quintic merger gives $S\mid T$ and the
coprime square gives $b$ square and $O>0$.  Sun's reverse bridge gives
a nonnegative zero of $P_0$.

Combining this with the fixed parent/compiler endpoint yields
\[
 \mathrm{ZFC}\vdash
 \left[\neg\exists(f,g,h,k,l,w,x,y,z,m)\in\mathbb Z^{10}\ F_{10}=0\right]
 \longleftrightarrow\operatorname{Con}(\mathrm{ZFC}).
\]
This invokes the stated mathematical theorems, including Sun's 2026
preprint.  The computational receipt verifies finite coefficient and
circuit claims; it is not an end-to-end formal proof of this equivalence.

\subsection{Reproduction and quantitative comparison}

The two commands \texttt{python build\_signed10.py} and
\texttt{python verify\_signed10.py}, run in the \texttt{signed10} directory,
use only the Python standard library.  The source sum is evaluated by a
shared-power Horner circuit; the geometric sum uses repeated doubling,
without division.  All constant powers are expanded into binary gates.
The main circuit has 566 additions, 25 subtractions and 851 multiplications,
with 79 fixed leaves.  A nonzero leading coefficient after diagonal
specialization modulo 101, 1009 and 1000003 proves that the circuit degree
upper bound is attained.

For transparency, the main degree calculation is as follows.  Here
$t=5^{31}$, and $X_n,A_i$ refer to the finalizer rather than the source
polynomial $X$:
\begin{center}
\begin{tabular}{@{}lr@{}}
\toprule
Expression & Total degree\\
\midrule
$X$ & $120t+24$\\
$Y$ & $80t+16$\\
$S$ & $280t+59$\\
$r_0$ & $1132t+250$\\
$r_1=R_2$ & $1920t+402$\\
$r_2=DFI^*$ & $24640t+5268$\\
$\Sigma$ & $1400t+297$\\
$X_n$ & $1400t+310$\\
$A_0$ & $3932t+844$\\
$A_1$ & $6984t+1496$\\
$A_2$ & $31624t+6766$\\
\bottomrule
\end{tabular}
\end{center}
The degree of $A_2$ strictly exceeds those of $X_n^2,A_0,A_1$.
Consequently the top part of the final norm is $A_2^4$, giving
$\deg F_{10}=126496t+27064$.  The recorded nonzero diagonal coefficients
are exact finite certificates that these top parts do not cancel for
the actual fixed coefficient tuple.

\begin{center}
\begin{tabular}{lrr}
Variant & Operations & Exact degree, $t=5^{31}$\\ \hline
Sun's four-square sextic specialization &1492&$2407680t+514640$\\
Four squares, quintic merger &1492&$2403200t+513696$\\
Three squares, Sun's generic weights &1445&$807360t+172560$\\
Three squares, universal successive weights &1444&$158592t+33880$\\
Three squares, structural signed weights &1442&$126496t+27064$
\end{tabular}
\end{center}
All five variants have the same ten unrestricted integer unknowns and
the same fixed source coefficients.  The complete proof, exact circuits,
coefficient recipes, binary payloads, source references and verification
receipts are in the \texttt{signed10} directory.
